\documentclass[11pt]{article}
% cs250preamble.tex --- shared preamble for CS 250 course notes
%
% This file is \input by main.tex and by each individual module
% file so that each module can still compile standalone. Any
% package or custom-environment change should be made here, not
% in the individual module files.

\usepackage[T1]{fontenc}
\usepackage[margin=1in]{geometry}
\usepackage{amsmath, amssymb, amsthm}
\usepackage{enumitem}
\usepackage{hyperref}
\usepackage{booktabs}
\usepackage{listings}
\usepackage{xcolor}
\usepackage{tikz}
\usetikzlibrary{shapes.geometric, shapes.gates.logic.US, shapes.gates.logic.IEC, arrows.meta, positioning, calc, trees}
\usepackage{bussproofs}
\usepackage{mdframed}

\lstset{
  basicstyle=\ttfamily\small,
  frame=single,
  breaklines=true,
  columns=fullflexible,
  mathescape=true
}

\theoremstyle{definition}
\newtheorem{theorem}{Theorem}[section]
\newtheorem{definition}[theorem]{Definition}
\newtheorem{example}[theorem]{Example}

\hypersetup{
  colorlinks=true,
  linkcolor=blue!60!black,
  urlcolor=blue!60!black
}

% Key-takeaway box: used at the end of major sections and at the end of
% each module to summarise the main points. Expects \item bullets inside.
\newmdenv[
  linewidth=0.8pt,
  linecolor=blue!40!black,
  backgroundcolor=blue!3,
  skipabove=8pt,
  skipbelow=8pt,
  innertopmargin=7pt,
  innerbottommargin=7pt,
  innerleftmargin=9pt,
  innerrightmargin=9pt
]{keytakeawaybox}
\newenvironment{keytakeaway}{%
  \begin{keytakeawaybox}%
  \noindent\textbf{Key takeaways.}%
  \begin{itemize}[leftmargin=1.3em, topsep=2pt, itemsep=1pt, label=\textbullet]%
}{%
  \end{itemize}%
  \end{keytakeawaybox}%
}

% Summary/looking-ahead environment: used once at the end of each module
% to wrap up what the module built and point forward.
\newmdenv[
  linewidth=0.8pt,
  linecolor=gray!60,
  backgroundcolor=gray!5,
  skipabove=8pt,
  skipbelow=8pt,
  innertopmargin=7pt,
  innerbottommargin=7pt,
  innerleftmargin=9pt,
  innerrightmargin=9pt
]{summarybox}

% Sidebar environment: used for tangential asides---historical notes,
% pointers to other areas of math/CS, cross-paradigm remarks---that
% enrich the main narrative without being essential to it. Takes one
% argument: the sidebar's title.
\newmdenv[
  linewidth=0.6pt,
  linecolor=gray!60,
  backgroundcolor=gray!4,
  skipabove=8pt,
  skipbelow=8pt,
  innertopmargin=6pt,
  innerbottommargin=6pt,
  innerleftmargin=9pt,
  innerrightmargin=9pt
]{sidebarbox}
\newenvironment{sidebar}[1]{%
  \begin{sidebarbox}%
  \noindent\textbf{Sidebar: #1.}\par\medskip%
}{%
  \end{sidebarbox}%
}

% Reading-map environment: used at the top of each numbered module to
% indicate Epp coverage, prerequisites, relevant intermezzi, and
% suggested practice problems. Expects four \rmentry{label}{body}
% items inside (or arbitrary paragraphs).
\newmdenv[
  linewidth=0.8pt,
  linecolor=black!55,
  backgroundcolor=yellow!4,
  skipabove=10pt,
  skipbelow=14pt,
  innertopmargin=9pt,
  innerbottommargin=9pt,
  innerleftmargin=11pt,
  innerrightmargin=11pt
]{readingmapbox}
\newcommand{\rmentry}[2]{\par\noindent\textbf{#1}\hspace{0.4em}#2\par}
\newenvironment{readingmap}{%
  \begin{readingmapbox}%
  \noindent{\bfseries\large Reading map}%
  \vspace{2pt}%
  \setlength{\parskip}{4pt plus 1pt}%
}{%
  \end{readingmapbox}%
}


\title{CS 250 --- Further Reading}
\author{}
\date{}

\begin{document}
\maketitle

Throughout this book, most intermezzi ended with a ``Further reading'' section
pointing at accessible next steps if the topic caught your interest. This
appendix collects those entries in one place, organized by topic, for readers
who want the whole reading list at a glance. Each entry is reproduced from the
intermezzo where it originally appeared, so there is some overlap when a
single source is relevant to multiple topics.

A few of the entries below are genuinely foundational papers that everyone in
computer science should read at some point; others are gentle on-ramps for
specific subtopics. If you only have time for one or two, the entries marked
\textbf{[start here]} are the ones to pick up first.

\section*{Natural Deduction and Proof Theory}
\emph{Related intermezzo: Inference Rules as Trees (Module~3).}

\begin{itemize}
  \item \textbf{[start here]} Frank Pfenning, \href{https://www.cs.cmu.edu/~fp/courses/15317-s23/lectures/02-natded.pdf}{\emph{Natural Deduction}} (lecture notes for CMU 15-317, ``Constructive Logic''). The single best gentle introduction to natural deduction in tree form, with exactly the discipline of introduction and elimination rules used in the intermezzo.
  \item Jan von Plato, ``\href{https://projecteuclid.org/journals/bulletin-of-symbolic-logic/volume-18/issue-3/Gentzens-proof-systems-byproducts-in-a-work-of-genius/10.2178/bsl/1344861886.short}{Gentzen's Proof Systems: Byproducts in a Work of Genius},'' \emph{Bulletin of Symbolic Logic} \textbf{18}(3), 313--367 (2012). A historical and conceptual essay on how Gentzen actually invented natural deduction and the sequent calculus, written for a logic-curious but non-specialist audience.
  \item Dag Prawitz, \emph{\href{https://store.doverpublications.com/0486446557.html}{Natural Deduction: A Proof-Theoretical Study}}, Almqvist \& Wiksell (1965); Dover reprint (2006). The first chapter is the canonical exposition of normalization---``every detour can be removed''---which is the proof-theory version of evaluating a program. Stop at the end of Chapter~I; later chapters are more technical.
\end{itemize}

\section*{Cantor, Infinity, and the Diagonal Argument}
\emph{Related intermezzo: Cardinality and Function Spaces (Module~1).}

\begin{itemize}
  \item Robert Gray, ``Georg Cantor and Transcendental Numbers,'' \emph{American Mathematical Monthly} \textbf{101}(9), 819--832 (1994). A clean historical account of Cantor's \emph{original} 1874 uncountability argument (which was not the diagonal argument!) and how the diagonal version came later. Excellent for seeing how the ideas in the intermezzo actually developed.
  \item \textbf{[start here]} Martin Davis, ``\href{https://archive.org/details/mathematicstoday00stee}{What is a Computation?}'' in L.~A.\ Steen (ed.), \emph{Mathematics Today: Twelve Informal Essays}, Springer (1978). Davis walks you from Cantor's diagonal argument straight to the unsolvability of the halting problem in a way that requires almost no prior CS background---perfect if you want to see the forward look from this chapter cashed out.
  \item Raymond Smullyan, \emph{\href{https://academic.oup.com/book/54024}{Diagonalization and Self-Reference}}, Oxford Logic Guides 27, Oxford University Press (1994). A whole book---but the early chapters are gentle, and they show how Cantor, G\"odel, and Tarski are all variations on a single self-referential theme.
\end{itemize}

\section*{Curry--Howard and Proofs as Programs}
\emph{Related intermezzi: Proofs Are Programs (Module~7), The Lambda Calculus (end of book).}

\begin{itemize}
  \item \textbf{[start here]} Philip Wadler, ``\href{https://homepages.inf.ed.ac.uk/wadler/papers/propositions-as-types/propositions-as-types.pdf}{Propositions as Types},'' \emph{Communications of the ACM} \textbf{58}(12), 75--84 (2015). \emph{The} accessible undergraduate-friendly tour of the Curry--Howard correspondence, with historical context, examples in modern languages, and Wadler's characteristic clarity. If you read only one paper from this list, read this one.
  \item William A.\ Howard, ``\href{https://www.cs.cmu.edu/~crary/819-f09/Howard80.pdf}{The Formulae-as-Types Notion of Construction},'' in J.\ R.\ Hindley and J.\ P.\ Seldin (eds.), \emph{To H.\ B.\ Curry: Essays on Combinatory Logic, Lambda Calculus and Formalism}, Academic Press (1980); manuscript circulated 1969. The original short note that started everything. Only about a dozen pages, and the first half is genuinely readable once you have seen natural deduction.
  \item Philip Wadler, ``\href{https://homepages.inf.ed.ac.uk/wadler/papers/frege/frege.pdf}{Proofs are Programs: 19th Century Logic and 21st Century Computing},'' \emph{Dr.\ Dobb's Journal} (2000). A shorter, more popular precursor to ``Propositions as Types,'' pitched at working programmers---good as a warm-up before tackling the CACM piece.
\end{itemize}

\section*{The Lambda Calculus}
\emph{Related intermezzo: The Lambda Calculus (end of book).}

\begin{itemize}
  \item Henk Barendregt, ``\href{https://www.cambridge.org/core/journals/bulletin-of-symbolic-logic/article/abs/impact-of-the-lambda-calculus-in-logic-and-computer-science/7B63F923EC566E549B093911BFC7470E}{The Impact of the Lambda Calculus in Logic and Computer Science},'' \emph{Bulletin of Symbolic Logic} \textbf{3}(2), 181--215 (1997). The standard expository overview of why the lambda calculus matters across logic and computer science, written by the field's leading authority. The first half is fully undergraduate-accessible.
  \item \textbf{[start here]} Peter Selinger, \emph{\href{https://arxiv.org/abs/0804.3434}{Lecture Notes on the Lambda Calculus}}, arXiv:0804.3434 (revised 2013). The cleanest free introduction I know of---covers Church encodings, beta reduction, the Church--Rosser theorem, and fixed-point combinators at exactly undergraduate pitch.
  \item Felice Cardone and J.\ Roger Hindley, ``\href{http://www.di.unito.it/~felice/pdf/lambdacomb.pdf}{History of Lambda-calculus and Combinatory Logic},'' in D.\ Gabbay and J.\ Woods (eds.), \emph{Handbook of the History of Logic}, vol.~5, Elsevier (2009). A historical narrative of how Church, Curry, and Turing arrived at the lambda calculus and combinatory logic---gives the Church--Turing-thesis story in human terms.
\end{itemize}

\section*{The Algebra of Types}
\emph{Related intermezzo: The Algebra of Types (end of book).}

\begin{itemize}
  \item \textbf{[start here]} Conor McBride, ``The Derivative of a
  Regular Type is its Type of One-Hole Contexts,'' unpublished
  manuscript (2001), available on the author's homepage at
  \url{https://strictlypositive.org}. A short, witty, surprisingly
  elementary derivation of why \emph{differentiating} an algebraic data
  type literally gives you its zipper. Connects high-school calculus to
  programming in a way undergraduates love.
  \item Brent Yorgey, ``Species and Functors and Types, Oh My!'' \emph{Haskell Symposium 2010}. Written explicitly as an accessible bridge between Joyal's combinatorial species and the algebraic data types programmers already know. The most undergraduate-friendly entry point to species I know of.
  \item Philippe Flajolet and Robert Sedgewick, \emph{\href{https://ac.cs.princeton.edu/home/}{Analytic Combinatorics}}, Cambridge University Press (2009), Chapter~I (``Symbolic Methods''). The full book PDF is freely available on Sedgewick's Princeton homepage. Chapter~I introduces the symbolic method---sums, products, sequences as type combinators with generating functions---at exactly the level a strong undergrad can handle.
\end{itemize}

\section*{Temporal Logic and Model Checking}
\emph{Related intermezzo: Other Logics, Other Worlds (Module~7).}

\begin{itemize}
  \item \textbf{[start here]} Chris Newcombe, Tim Rabbat, Fan Zhang, Bogdan Munteanu, Marc Brooker, and Michael Deardeuff, ``How Amazon Web Services Uses Formal Methods,'' \emph{Communications of the ACM} \textbf{58}(4), 66--73 (2015). A wonderfully accessible report by practicing engineers on how AWS uses TLA+ to verify the correctness of distributed systems like S3 and DynamoDB. If you wonder why anyone outside academia cares about temporal logic, this is your answer.
  \item Edmund M.\ Clarke, E.\ Allen Emerson, and Joseph Sifakis, ``Model Checking: Algorithmic Verification and Debugging,'' \emph{Communications of the ACM} \textbf{52}(11), 74--84 (2009). The Turing Award lecture: a retrospective on how model checking was developed and why it works. Written for a general CS audience.
  \item Amir Pnueli, ``The Temporal Logic of Programs,'' \emph{18th Annual Symposium on Foundations of Computer Science} (FOCS), 46--57 (1977). The founding paper of program temporal logic. Sections 1--2 are historically essential and undergraduate-readable; later sections get more technical.
\end{itemize}

\section*{Equivalence, Sameness, and Structuralism}
\emph{Related intermezzo: What Does It Mean for Things to Be ``The Same''? (Module~6).}

\begin{itemize}
  \item \textbf{[start here]} Barry Mazur, ``\href{https://people.math.osu.edu/cogdell.1/6112-Mazur-www.pdf}{When is one thing equal to some other thing?}'' in \emph{Proof and Other Dilemmas: Mathematics and Philosophy} (B.~Gold and R.~A.\ Simons, eds.), Mathematical Association of America (2008). The single best expository essay for undergraduates on equivalence, quotient, and ``sameness.'' Mazur writes beautifully and assumes almost nothing.
  \item Steve Awodey, ``\href{https://www.andrew.cmu.edu/user/awodey/preprints/siu.pdf}{Structuralism, Invariance, and Univalence},'' \emph{Philosophia Mathematica} \textbf{22}(1), 1--11 (2014). A short, accessible philosophical essay on what it means for two mathematical structures to be ``the same'' and on why equivalence (rather than literal equality) is the right notion. The closing section connects to the modern Univalence axiom in homotopy type theory.
  \item Paul Halmos, \emph{Naive Set Theory}, Springer (1960; reprint 1974), Sections 7--8. Not a paper, but Halmos's two-page treatment of equivalence relations and partitions is the canonical undergraduate reference and pairs nicely with the philosophical pieces above.
\end{itemize}

\section*{Game Semantics, Interactive Proofs, and Zero Knowledge}
\emph{Related intermezzo: Quantifiers as Games (Module~4).}

\begin{itemize}
  \item Jaakko Hintikka, ``Language-Games for Quantifiers,'' \emph{American Philosophical Quarterly} Monograph Series no.~2 (1968), 46--72. Hintikka's own short and accessible exposition of how quantifiers are best read as a two-player game between a Verifier and a Falsifier---the philosophical source of the framework in the intermezzo.
  \item Shafi Goldwasser, Silvio Micali, and Charles Rackoff, ``The Knowledge Complexity of Interactive Proof Systems,'' \emph{SIAM Journal on Computing} \textbf{18}(1), 186--208 (1989). The foundational paper on interactive proofs and zero-knowledge; widely available via Silvio Micali's MIT homepage and elsewhere. The first few sections, where they set up the Prover/Verifier game, are surprisingly accessible and put a precise mathematical frame around the picture sketched in the chapter.
  \item \textbf{[start here]} Oded Goldreich, \emph{\href{https://www.wisdom.weizmann.ac.il/~oded/zk-tut02.html}{Zero-Knowledge: A Tutorial}}, available on the author's homepage at the Weizmann Institute. Written explicitly to be the most accessible on-ramp to zero-knowledge proofs for non-specialists. If interactive proofs caught your attention, this is where to go next.
\end{itemize}

\section*{Induction and the Peano Axioms}
\emph{Related intermezzo: More on Natural Induction (Module~4).}

\begin{itemize}
  \item \textbf{[start here]} Leon Henkin, ``On Mathematical Induction,'' \emph{American Mathematical Monthly} \textbf{67}(4), 323--338 (1960). The classic expository article distinguishing the \emph{principle} of induction from the \emph{axiom}, and showing precisely which Peano axioms are needed to prove commutativity and associativity of addition---the same proofs the intermezzo walks through.
  \item George P\'olya, \emph{\href{https://press.princeton.edu/books/paperback/9780691025094/mathematics-and-plausible-reasoning-volume-1}{Induction and Analogy in Mathematics}} (volume~I of \emph{Mathematics and Plausible Reasoning}), Princeton University Press (1954), particularly the chapter on mathematical induction. P\'olya's gentle, example-driven exposition is aimed at students discovering proof for the first time.
  \item Edmund Landau, \emph{Foundations of Analysis}, Chelsea (1951; original German 1930), Chapter~1. Landau builds the natural numbers from the Peano axioms and proves commutativity and associativity of addition with full inductive rigor in about ten pages---you can literally check every step. Terse but transparent.
\end{itemize}

\section*{Modular Arithmetic and Finite Fields}
\emph{Related module: Module~6.}

\begin{itemize}
  \item \emph{Wikipedia: \href{https://en.wikipedia.org/wiki/Finite_field_arithmetic}{Finite field arithmetic}}. Useful starting reference for how finite fields are actually computed with in software and hardware.
  \item \emph{Wikipedia: \href{https://en.wikipedia.org/wiki/Finite_field\#Discrete_logarithm}{Finite field: Discrete logarithm}}. The discrete logarithm problem in a finite field is the hard problem underlying Diffie--Hellman key exchange and a surprising amount of modern cryptography.
\end{itemize}

\end{document}
